% ---------------------------------------------------------------------
% Parte 5 · Panel de Erwin · expone Erwin · 1:30 (lo edita solo Erwin)
% La captura es del panel funcionando en tu proyecto, tomada por ti y con
% la fecha visible: guárdala como figuras/panel-erwin.png.
% Las tres tarjetas tienen la misma estructura en los tres paneles para
% que la comparación de la parte 8 salga sola.
% ---------------------------------------------------------------------
\begin{frame}{Panel de Erwin}{\pend{Erwin}{qué proyecto es y para quién es el panel, en una línea}}
  \begin{columns}[T]
    \begin{column}{0.56\textwidth}
      \figura{figuras/panel-erwin.png}{\linewidth}{3.7cm}{Captura del panel funcionando, con la fecha
        visible (evidencia de primera mano).}
      \begin{tarjeta}{bcmAcento}{Si se cae la conexión}
        \footnotesize\pend{Erwin}{qué ve el usuario y cómo se recupera}
      \end{tarjeta}
    \end{column}
    \begin{column}{0.4\textwidth}
      \begin{tarjeta}{bcmTinta}{Qué hace}
        \footnotesize\pend{Erwin}{qué muestra o controla el panel}
      \end{tarjeta}
      \begin{tarjeta}{bcmWS}{Dónde entra el WebSocket}
        \footnotesize\pend{Erwin}{qué evento viaja, quién lo envía y quién lo recibe}
      \end{tarjeta}
    \end{column}
  \end{columns}
\end{frame}
